\section{Objectifs}
\begin{itemize}
\item Relier taux de variation, sécantes et nombre dérivé.
\item Utiliser le nombre dérivé comme pente de tangente et comme taux instantané.
\item Déterminer l'équation d'une tangente et une approximation linéaire.
\item Calculer les dérivées des fonctions usuelles et de sommes, produits, inverses et quotients.
\end{itemize}

\section{Cours}
\subsection{Point de vue local}
Le taux de variation entre \texttt{a} et \texttt{a+h} vaut \texttt{(f(a+h)-f(a))/h} pour \texttt{h≠0}. Graphiquement, c'est la pente d'une sécante. Lorsque ce taux se stabilise quand \texttt{h} devient proche de zéro, on introduit intuitivement le nombre dérivé \texttt{f'(a)}.

La tangente au point d'abscisse \texttt{a} a pour équation \texttt{y=f(a)+f'(a)(x-a)}. Pour \texttt{h} petit, \texttt{f(a+h)≈f(a)+f'(a)h} : c'est l'approximation linéaire.

Le nombre dérivé peut représenter une vitesse instantanée, un coût marginal ou un autre taux instantané.

\subsection{Point de vue global}
Une fonction dérivable sur un intervalle possède une fonction dérivée. Formules de base : \texttt{(x²)'=2x}, \texttt{(x³)'=3x²}, \texttt{(1/x)'=-1/x²}, \texttt{(√x)'=1/(2√x)} pour \texttt{x>0}, et plus généralement \texttt{(x^n)'=n x^{n-1}} dans le domaine adapté pour \texttt{n∈Z}.

Règles : \texttt{(u+v)'=u'+v'}, \texttt{(uv)'=u'v+uv'}, \texttt{(1/u)'=-u'/u²}, \texttt{(u/v)'=(u'v-uv')/v²} lorsque les dénominateurs ne s'annulent pas.

La fonction valeur absolue n'est pas dérivable en 0 ; la racine carrée n'est pas dérivable en 0.

\section{Démonstrations à travailler}
\begin{itemize}
\item équation de la tangente ;
\item non-dérivabilité de \texttt{√x} en 0 ;
\item dérivées de \texttt{x²} et \texttt{1/x} ;
\item dérivée d'un produit.
\end{itemize}

\section{Erreurs fréquentes}
\begin{itemize}
\item Confondre \texttt{f'(a)} avec \texttt{f(a)}.
\item Oublier le domaine lors d'une dérivée d'inverse, quotient ou racine.
\item Traiter la limite du taux de variation avec un formalisme de Terminale non requis ici.
\end{itemize}

\section{Synthèse}
La dérivation relie une variation locale à une pente. Le nombre dérivé décrit un comportement en un point ; la fonction dérivée permet ensuite d'étudier globalement une fonction.
